To obtain a single maximizer, it suffices to control one maximizing sequence. The hypothesis controls only the chosen sequence, with no uniform bound over the entire current family.

\begin{proposition}\label{prop:selected-gauduchon-attainment}
Let $0\leq p\in\mathcal D_\omega$ be nonzero and put $R=\rho_\omega(p)$. Suppose $T_j\in\mathscr G_\omega$ converge weakly to $T$ and satisfy $D_p(T_j)\to R$. If the densities $p k_{T_j}$ are uniformly integrable, then
\[
 D_p(T)=R,\qquad
 p k_{T_j}\,\mathrm{d}V\rightharpoonup p k_T\,\mathrm{d}V.
\]
Consequently, a uniformly integrable maximizing sequence has a maximizing current among its weak cluster limits. No assumption $p\in\mathcal D_\omega^0$ is needed.
\end{proposition}
\begin{proof}
Put $p_K=\min(p,K)$. Since $p$ is finite almost everywhere, uniform integrability gives
\[
 \varepsilon_K:=\sup_j\int_{\{p>K\}}p k_{T_j}\,\mathrm{d}V
       \longrightarrow0.
\]
For every nonnegative continuous $h$, the weight $hp_K$ is bounded and measurable. The bounded-weight upper semicontinuity proved in \cref{thm:integrable-current-attainment} therefore gives
\[
 \limsup_j\int_Xhp k_{T_j}\,\mathrm{d}V
 \leq \int_Xhp_K k_T\,\mathrm{d}V
                   +\|h\|_\infty\varepsilon_K.
\]
Let $K\to\infty$. For $h=1$ this yields $R\leq D_p(T)$; the supremum formula \eqref{eq:integrable-current-supremum} gives the reverse inequality. If $\mu$ is a weak cluster limit of $p k_{T_j}\,\mathrm{d}V$, the same estimate for all nonnegative continuous $h$ gives $\mu\leq p k_T\,\mathrm{d}V$. Both measures have mass $R$, so they agree. Their bounded masses give weak compactness, and every cluster limit has now been identified. The final assertion follows from compactness of $\mathscr G_\omega$.
\end{proof}

This criterion gives convergence of determinant-weighted measures at the chosen maximum. It supplies neither smoothness of the current nor an exact Monge--Amp\`ere equation.
